% This appendix is a manually curated publication view.
% It intentionally omits repository paths, file names, scripts, and commands.

\chapter{实验参数、评价方法与证据边界}

本附录汇总正文中不宜反复展开、但对学术复核仍有必要的实验设定、数学定义与证据边界。内容按场景设计、故障模型、控制参数、优化约束、评价指标和证据等级组织。变量均在首次出现时说明其意义；仅用于定位代码或执行实验的路径、文件名、脚本名和命令不属于论文论证链，故不在本附录保留。

除特别说明外，表中数值是本文仿真与数字孪生实验的名义设定，不代表传感器制造误差、海试环境统计量或装备验收限值。凡尚未由实物数据标定的参数，均只能用于可控扰动下的相对比较。

\section{仿真场景与扰动因子}\label{sec:app-a-1}

场景库采用“单因素分级 + 多因素联合”的设计。单因素场景用于识别某类观测退化对估计与控制的独立影响，综合压力场景用于检查多种不确定性同时出现时的闭环鲁棒性。令 \(p_{\mathrm{dvl}}\) 表示 DVL 观测在给定时段内的丢失概率，\(B_{\mathrm{sat}}\) 表示磁场硬饱和阈值，\(\eta_a\) 与 \(\eta_s\) 分别表示加速度计和声呐噪声的相对放大系数，\(v_c\) 表示海流速度。

\begin{longtable}{@{}p{0.22\linewidth}p{0.20\linewidth}p{0.20\linewidth}p{0.28\linewidth}@{}}
\caption{仿真场景及其学术用途}\label{tab:app-a-scenarios}\\
\toprule
场景 & 主要扰动 & 名义强度 & 用途 \\
\midrule
\endfirsthead
\toprule
场景 & 主要扰动 & 名义强度 & 用途 \\
\midrule
\endhead
\bottomrule
\endfoot
干净基线 & 无主动扰动 & --- & 给出算法对照基线 \\
DVL 轻度丢失 & 速度观测间歇缺失 & \(p_{\mathrm{dvl}}=0.10\) & 检查轻度观测不连续 \\
DVL 中度丢失 & 速度观测间歇缺失 & \(p_{\mathrm{dvl}}=0.30\) & 检查常见遮挡压力 \\
DVL 重度丢失 & 速度观测间歇缺失 & \(p_{\mathrm{dvl}}=0.60\) & 检查置信度调度边界 \\
DVL 极重度丢失 & 速度观测近乎不可用 & \(p_{\mathrm{dvl}}=0.90\) & 检查估计器失效趋势 \\
轻度磁畸变 & 磁场限幅 & \(B_{\mathrm{sat}}=10^{-6}\,\mathrm{T}\) & 检查磁观测降权机制 \\
重度磁畸变 & 磁场限幅 & \(B_{\mathrm{sat}}=10^{-8}\,\mathrm{T}\) & 检查磁观测近失效状态 \\
声呐杂波 & 成像观测噪声增强 & 固定放大 & 检查声学质量退化 \\
综合压力 & 丢包、漂移、噪声与海流联合 & 见表~\ref{tab:app-a-combined-stress} & 检查多扰动叠加鲁棒性 \\
\end{longtable}

综合压力场景同时施加八类扰动。令 \(p_b\) 表示通信帧丢失概率，\(\dot b_a\) 表示加速度计偏置增长率，\(\lambda_z\) 表示深度尖峰的平均发生率，\(A_z\) 表示尖峰幅值。各量的名义值与物理意图列于表~\ref{tab:app-a-combined-stress}。

\begin{longtable}{@{}p{0.24\linewidth}p{0.28\linewidth}p{0.38\linewidth}@{}}
\caption{综合压力场景的扰动参数}\label{tab:app-a-combined-stress}\\
\toprule
扰动源 & 名义参数 & 物理意图 \\
\midrule
\endfirsthead
\toprule
扰动源 & 名义参数 & 物理意图 \\
\midrule
\endhead
\bottomrule
\endfoot
通信帧丢失 & \(p_b=0.05\) & 近似低带宽链路中的随机丢帧 \\
DVL 丢失 & \(p_{\mathrm{dvl}}=0.30\)，持续 \(1\sim3\,\mathrm{s}\) & 近似浑水或遮挡引起的速度观测中断 \\
IMU 慢漂移 & \(\dot b_a=10^{-3}\,\mathrm{m/s^3}\) & 近似无充分温补时的缓慢偏置增长 \\
深度尖峰 & \(\lambda_z=0.05\,\mathrm{Hz}\)，\(A_z=0.5\,\mathrm{m}\) & 近似低频瞬态电磁干扰 \\
磁场饱和 & \(B_{\mathrm{sat}}=10^{-7}\,\mathrm{T}\) & 近似强近场磁干扰导致的量程受限 \\
加速度计噪声 & \(\eta_a=1.5\) & 近似结构振动耦合 \\
声呐噪声 & \(\eta_s=2.0\) & 近似浑浊浅水中的成像退化 \\
海流 & \(v_c=0.3\,\mathrm{m/s}\) & 近似常见近海横流压力 \\
\end{longtable}

上述磁饱和阈值只用于形成可重复的相对压力，不是由实测噪声谱反演得到的传感器参数。因此，本文只比较同一场景定义下不同算法的相对行为，不把该阈值解释为真实海域中的失效概率。

\section{异步采样与故障注入模型}\label{sec:app-a-2}

仿真接口由通信时延、异步采样和故障注入三个抽象环节组成。通信时延写为
\[
  \tau_k=\bar{\tau}+\epsilon_{\tau,k},
  \qquad
  \bar{\tau}=200\,\mathrm{ms},
  \qquad
  \epsilon_{\tau,k}\sim\mathcal{U}(-50,50)\,\mathrm{ms},
\]
其中 \(\bar{\tau}\) 为基线时延，\(\epsilon_{\tau,k}\) 为有界抖动。各传感器按独立时钟采样：IMU、深度、磁场和 DVL 的名义频率分别为 \(100\)、\(50\)、\(20\) 和 \(6\,\mathrm{Hz}\)。该设定用于复现多速率观测的时间对齐问题，不代表所有实物传感器的固定输出频率。

六类故障模型定义如下。

\begin{longtable}{@{}p{0.20\linewidth}p{0.35\linewidth}p{0.35\linewidth}@{}}
\caption{故障注入模型及其参数}\label{tab:app-a-chaos}\\
\toprule
故障类型 & 数学或逻辑模型 & 名义参数 \\
\midrule
\endfirsthead
\toprule
故障类型 & 数学或逻辑模型 & 名义参数 \\
\midrule
\endhead
\bottomrule
\endfoot
DVL 冻结 & 丢失窗口内保持最近一次有效速度 \(\tilde{\boldsymbol v}_k=\tilde{\boldsymbol v}_{k_0}\) & 触发概率 \(0.30\)，窗口 \(1\sim3\,\mathrm{s}\) \\
IMU 漂移 & \(\tilde{\boldsymbol a}(t)=\boldsymbol a(t)+\boldsymbol b_a(t)\)，\(\dot{\boldsymbol b}_a=\boldsymbol w_b\) & 名义漂移率 \(10^{-3}\,\mathrm{m/s^3}\) \\
深度尖峰 & \(\tilde z_k=z_k+A_z I_k\)，\(I_k\) 服从泊松到达过程 & \(\lambda_z=0.05\,\mathrm{Hz}\)，\(A_z=0.5\,\mathrm{m}\) \\
磁场饱和 & \(\tilde B_i=\operatorname{clip}(B_i,-B_{\mathrm{sat}},B_{\mathrm{sat}})\) & \(B_{\mathrm{sat}}=10^{-7}\,\mathrm{T}\) \\
通信丢包 & 每帧按伯努利变量 \(I_k\sim\operatorname{Bernoulli}(1-p_b)\) 保留或丢弃 & \(p_b=0.05\) \\
数据乱序 & 以给定概率交换相邻帧的到达顺序 & 乱序概率 \(0.01\) \\
\end{longtable}

这里的 \(\operatorname{clip}(\cdot)\) 表示双边限幅，\(\boldsymbol w_b\) 表示驱动偏置变化的随机过程。模型分别覆盖“持续值错误、缓慢累积误差、瞬态异常、量程受限和通信时序错误”，便于区分估计器对不同失效形态的响应。

\section{不确定性感知控制参数化与消融设计}\label{sec:app-a-3}

令 \(c\in[0,1]\) 为由状态协方差映射得到的归一化置信度，\(c=1\) 表示高置信度。定义低置信度激活量
\[
  \rho(c)=(1-c)^\alpha ,
\]
其中 \(\alpha>0\) 控制置信度映射的非线性。跟踪权重和控制增量权重分别写为
\[
  w_e(c)=w_{e,0}\left[1+s_e\rho(c)\right],
\]
\[
  w_u(c)=w_{u,0}
  \left[
    \gamma_u+(1-\gamma_u)
    \sigma\!\left(k(c-c_0)\right)
  \right],
  \qquad
  \sigma(x)=\frac{1}{1+\exp(-x)}.
\]
其中 \(w_{e,0}\) 与 \(w_{u,0}\) 为基准权重，\(s_e\) 为低置信度跟踪权重增益，\(\gamma_u\) 为控制代价的最小缩放比例，\(k\) 为 S 型函数陡峭度，\(c_0\) 为置信度转折点。

\begin{longtable}{@{}p{0.20\linewidth}p{0.17\linewidth}p{0.53\linewidth}@{}}
\caption{不确定性感知控制的名义参数}\label{tab:app-a-uampc-parameters}\\
\toprule
参数 & 名义值 & 含义 \\
\midrule
\endfirsthead
\toprule
参数 & 名义值 & 含义 \\
\midrule
\endhead
\bottomrule
\endfoot
\(s_e\) & \(3.0\) & 低置信度下跟踪权重的最大附加增益 \\
\(\gamma_u\) & \(0.3\) & 低置信度下控制增量代价的最小缩放比例 \\
\(\alpha\) & \(1.5\) & 置信度到跟踪增益的幂指数 \\
\(k\) & \(8.0\) & S 型控制代价缩放的陡峭度 \\
\(c_0\) & \(0.6\) & 高、低置信度过渡的中心值 \\
\(N\) & \(20\) & 模型预测控制的预测步数 \\
\(\Delta t\) & \(0.1\,\mathrm{s}\) & 离散控制时间步长 \\
\end{longtable}

消融实验不以工程开关命名，而按被移除的机制定义。

\begin{longtable}{@{}p{0.12\linewidth}p{0.25\linewidth}p{0.30\linewidth}p{0.23\linewidth}@{}}
\caption{不确定性感知控制消融变体}\label{tab:app-a-uampc-ablation}\\
\toprule
变体 & 定义 & 被隔离的机制 & 学术用途 \\
\midrule
\endfirsthead
\toprule
变体 & 定义 & 被隔离的机制 & 学术用途 \\
\midrule
\endhead
\bottomrule
\endfoot
A0 & \(w_e,w_u\) 恒定 & 全部置信度调度 & 固定权重对照 \\
A1 & 使用完整参数化 & 无 & 默认 UA-MPC \\
A2 & 令 \(w_u(c)=w_{u,0}\) & S 型控制代价降权 & 识别可解性与精度主贡献 \\
A3 & 令 \(\alpha=1\) & 非线性置信度映射 & 检查线性映射的边际影响 \\
A4 & 令 \(s_e=0\) & 低置信度跟踪增益 & 分离跟踪权重放大作用 \\
\end{longtable}

第 5 章的机制归因表明，A2 对可解性和 S 弯横向精度的影响最大；A3 只有边际收益。因此本文保留默认 \(\alpha=1.5\)，不依据单次运动学基准修改默认参数。

\section{代价权重与约束}\label{sec:app-a-4}

名义代价函数写为
\[
  J=
  \sum_{k=0}^{N-1}
  \left[
    \boldsymbol e_k^\mathsf{T}\boldsymbol Q(c)\boldsymbol e_k
    +\Delta\boldsymbol u_k^\mathsf{T}
     \boldsymbol R(c)\Delta\boldsymbol u_k
  \right]
  +\boldsymbol e_N^\mathsf{T}\boldsymbol Q_f\boldsymbol e_N ,
\]
其中 \(\boldsymbol e_k\) 为状态跟踪误差，\(\Delta\boldsymbol u_k\) 为控制增量，\(\boldsymbol Q(c)\) 与 \(\boldsymbol R(c)\) 分别为随置信度变化的跟踪和控制权重矩阵。

\begin{longtable}{@{}p{0.27\linewidth}p{0.18\linewidth}p{0.45\linewidth}@{}}
\caption{代价函数的基准权重}\label{tab:app-a-cost-weights}\\
\toprule
代价项 & 基准权重 & 说明 \\
\midrule
\endfirsthead
\toprule
代价项 & 基准权重 & 说明 \\
\midrule
\endhead
\bottomrule
\endfoot
纵向位置误差 \(e_x\) & \(1.0\) & 水平面位置跟踪 \\
横向位置误差 \(e_y\) & \(1.0\) & 电缆中心线跟踪 \\
深度误差 \(e_z\) & \(5.0\) & 直接关联近底安全裕度 \\
航向误差 \(e_\psi\) & \(3.0\) & 约束路径几何跟随 \\
前向速度误差 \(e_u\) & \(0.5\) & 允许适度速度调整 \\
垂向速度误差 \(e_w\) & \(1.0\) & 限制垂向动态偏差 \\
航向指令增量 & \(0.1\) & 抑制航向指令抖动 \\
深度指令增量 & \(0.1\) & 抑制深度指令跳变 \\
推力指令增量 & \(0.05\) & 抑制推进器频繁调节 \\
\end{longtable}

硬约束用于保证物理可行性，速率约束用于限制相邻预测步的变化，带宽约束用于限制指令相对当前状态的偏离。设 \(u_k\) 为前向速度，\(\psi_{\mathrm{cmd},k}\)、\(z_{\mathrm{cmd},k}\) 和 \(T_{\mathrm{cmd},k}\) 分别为航向、深度和推力指令。

\begin{longtable}{@{}p{0.27\linewidth}p{0.38\linewidth}p{0.25\linewidth}@{}}
\caption{模型预测控制约束}\label{tab:app-a-hard-constraints}\\
\toprule
约束 & 数学表达 & 名义范围 \\
\midrule
\endfirsthead
\toprule
约束 & 数学表达 & 名义范围 \\
\midrule
\endhead
\bottomrule
\endfoot
最低航速 & \(u_k\ge u_{\min}\) & \(u_{\min}=0.1\,\mathrm{m/s}\) \\
航向指令 & \(-\pi\le\psi_{\mathrm{cmd},k}\le\pi\) & \([-\pi,\pi]\,\mathrm{rad}\) \\
深度指令 & \(0\le z_{\mathrm{cmd},k}\le z_{\max}\) & \(z_{\max}=50\,\mathrm{m}\) \\
推力指令 & \(0\le T_{\mathrm{cmd},k}\le100\%\) & 禁止反向推力 \\
深度变化率 & \(\lvert\Delta z_{\mathrm{cmd},k}\rvert\le\Delta z_{\max}\) & \(\Delta z_{\max}=0.5\,\mathrm{m/step}\) \\
航向变化率 & \(\lvert\Delta\psi_{\mathrm{cmd},k}\rvert\le\Delta\psi_{\max}\) & \(\Delta\psi_{\max}=8^\circ/\mathrm{step}\) \\
深度带宽 & \(\lvert z_{\mathrm{cmd},k}-z_{\mathrm{current}}\rvert\le z_b\) & \(z_b=3.0\,\mathrm{m}\) \\
航向带宽 & \(\lvert\psi_{\mathrm{cmd},k}-\psi_{\mathrm{current}}\rvert\le\psi_b\) & \(\psi_b=45^\circ\) \\
\end{longtable}

其中 \(\Delta z_{\mathrm{cmd},k}=z_{\mathrm{cmd},k+1}-z_{\mathrm{cmd},k}\)，航向差采用角度环绕后的最短差值。硬约束、速率约束和带宽约束分别对应“物理边界、短时平滑性和中期可跟随性”，不能相互替代。

\section{数值求解、实时性与回退}\label{sec:app-a-5}

优化问题采用内点法求解。每个控制周期只更新当前状态、参考轨迹和置信度，并以上一周期最优控制序列移位后作为初值。名义配置见表~\ref{tab:app-a-solver}。

\begin{longtable}{@{}p{0.25\linewidth}p{0.22\linewidth}p{0.43\linewidth}@{}}
\caption{数值求解的名义配置}\label{tab:app-a-solver}\\
\toprule
项目 & 名义值 & 说明 \\
\midrule
\endfirsthead
\toprule
项目 & 名义值 & 说明 \\
\midrule
\endhead
\bottomrule
\endfoot
收敛容差 & \(10^{-4}\) & 平衡数值精度与在线计算时间 \\
最大迭代次数 & \(100\) & 防止单次求解无限阻塞控制循环 \\
控制周期 & \(100\,\mathrm{ms}\) & 与 \(\Delta t=0.1\,\mathrm{s}\) 一致 \\
预测时域 & \(N=20\) & 对应约 \(2\,\mathrm{s}\) 的前视窗口 \\
优化变量规模 & 约 \(360\) & 由状态维度、控制维度和预测步数共同决定 \\
\end{longtable}

若求解器在当前周期未返回可接受解，则系统显式进入回退逻辑，而不把回退输出记为求解成功。回退顺序为：优先保持最近一次已验证的安全控制输出；若该输出已过期，则退化到几何制导或安全设定值；当状态新鲜度或安全条件不满足时，由更高优先级的行为树与底层保护接管。

三场景、三种子扫描中，最大迭代次数为 100、200 和 400 时的事件加权回退率分别为 0.699、0.685 和 0.594，当前可观测失败占用分别为 88.41、192.21 和 372.20 ms。400 档的诊断帧数较 100 档减少约 26\%，表明阻塞已经降低控制更新频率。因此迭代上限不是本文宣称的性能优化手段，后续应优先改善初值连续性、约束可行性和终端任务定义。

\section{评价指标与数据契约}\label{sec:app-a-6}

本节只保留与结论直接相关的指标定义，不记录分析工具名称或数据通道名称。设第 \(k\) 个样本的真实位置和估计位置分别为 \(\boldsymbol p_k\) 与 \(\hat{\boldsymbol p}_k\)，水平误差为
\[
  e_{xy,k}=
  \sqrt{(p_{x,k}-\hat p_{x,k})^2+(p_{y,k}-\hat p_{y,k})^2}.
\]
水平位置均方根误差定义为
\[
  \operatorname{RMSE}_{xy}
  =
  \sqrt{\frac{1}{n}\sum_{k=1}^{n}e_{xy,k}^2}.
\]
\(\operatorname{CEP}_{50}\) 为样本集合 \(\{e_{xy,k}\}\) 的中位数，即满足半数水平误差不超过该值的半径。

对新息向量 \(\boldsymbol y_k\) 和新息协方差 \(\boldsymbol S_k\)，归一化新息平方定义为
\[
  \operatorname{NIS}_k=
  \boldsymbol y_k^\mathsf{T}
  \boldsymbol S_k^{-1}
  \boldsymbol y_k .
\]
若量测维数为 \(m\) 且滤波假设成立，则 \(\operatorname{NIS}_k\) 近似服从自由度为 \(m\) 的卡方分布，其期望为 \(m\)。跨不同量测维数比较时，另采用维数归一化量 \(\overline{\operatorname{NIS}}_k=\operatorname{NIS}_k/m\)，其理论期望为 \(1\)。本文使用 NIS 检查协方差一致性，不把它解释为在线定位精度本身。

控制量变化率采用相邻控制向量差分的均方根：
\[
  r_{\Delta u}
  =
  \sqrt{
    \frac{1}{n-1}
    \sum_{k=1}^{n-1}
    \lVert\boldsymbol u_{k+1}-\boldsymbol u_k\rVert_2^2
  }.
\]
求解回退率定义为
\[
  r_{\mathrm{fb}}=
  \frac{N_{\mathrm{fb}}}{N_{\mathrm{ctrl}}},
\]
其中 \(N_{\mathrm{fb}}\) 为进入回退逻辑的控制周期数，\(N_{\mathrm{ctrl}}\) 为总控制周期数。

对地形跟随，令 \(h_k\) 为真实离底高度，\(h^\star\) 为目标离底高度，则
\[
  \operatorname{RMSE}_{h}
  =
  \sqrt{\frac{1}{n}\sum_{k=1}^{n}(h_k-h^\star)^2},
  \qquad
  h_{\min}=\min_k h_k .
\]
对海缆巡检，令 \(d_{\perp,k}\) 为相对电缆中心线的横向偏移，\(s_k\) 为沿先验路由的累计进度，\(\hat d_k\) 与 \(\sigma_{d,k}\) 分别为埋深估计及其标准差。上述变量共同构成巡检质量与验收判断的最小数据契约。

双脑通信的学术数据契约只要求下行数据包含任务模式、控制意图和执行参数，上行数据包含姿态、速度、深度、能源状态与故障状态。具体字段偏移和程序接口属于实现文档，不进入论文附录；实物阶段仍需单独报告端到端时延分位数、丢包率和序列号跳变统计。

\section{证据矩阵与样本量边界}\label{sec:app-a-7}

原始数据存储位置会随平台和实验批次变化，不构成学术结论的一部分。本节以实验对象、样本量、可支撑结论和不可外推范围替代路径索引。

\begingroup
\small
\setlength{\tabcolsep}{3pt}
\renewcommand{\arraystretch}{1.04}
\begin{longtable}{@{}p{0.235\linewidth}p{0.185\linewidth}p{0.285\linewidth}p{0.235\linewidth}@{}}
\caption{核心实验的证据矩阵}\label{tab:app-a-data-sources}\\
\toprule
实验对象 & 样本量 & 可支撑结论 & 不可外推范围 \\
\midrule
\endfirsthead
\toprule
实验对象 & 样本量 & 可支撑结论 & 不可外推范围 \\
\midrule
\endhead
\bottomrule
\endfoot
状态估计扰动回归 & 8 场景 \(\times\) 3 种子 & 24 次均完成，输出稳定 & 不证明误差随扰动强度单调变化 \\
协方差一致性 & 8 场景 \(\times\) 3 种子 & 自适应机制在全部场景被触发 & NIS 为离线一致性检查 \\
UA-MPC 主消融 & 3 场景 \(\times\) 2 模式 \(\times\) 3 种子 & 综合压力下定位和横向指标同向改善 & 重度 DVL 丢失下不再有优势 \\
UA-MPC 可解性归因 & P-1 三变体 27 次；A1 默认层 9 次 & S 型控制代价降权为可解性主贡献 & 失败耗时尚未覆盖全部回退类型 \\
最大迭代次数对照 & 3 档 \(\times\) 3 场景 \(\times\) 3 种子 & 单纯增加迭代上限不能工程化修复高回退 & PVS/Mock AMD 结果不等同真机实时性 \\
横向机动归因 & 2 路径 \(\times\) 4 变体，确定性单次 & S 型控制代价降权为精度主贡献 & 不替代六自由度多种子结论 \\
地形跟随复验 & 3 强度 \(\times\) 3 种子 & 无安全违规，PID 为可靠基线 & 不证明 MPC 在所有地形下更优 \\
复杂平面路径 & 每路径单次 & MPC 在不规则曲率路径具备预瞄优势 & 直角路径上 LOS 仍可能更优 \\
TMR8637 随附测试报告 & 1 份报告，编号 6\# 模组三轴 & 1 Hz 本底噪声为 200 pT/\(\sqrt{\mathrm{Hz}}\)，50 Hz 约为 20--30 pT/\(\sqrt{\mathrm{Hz}}\) & 供应方随附单件报告，不等同于完整采集链实测 \\
磁外参仿真标定 & 单次全流程 & 标定计算与应用链路成立 & 不等同于真机转台或九参数标定 \\
代理电缆场景 & 6 场景 \(\times\) 2 模式，单种子 & 压力测试链路可运行 & 不能比较两种控制模式的统计优劣 \\
数字孪生验收 & 3 次独立运行 & 窗口化评分与聚合判据可复现 & 不等同于现场海试验收 \\
六自由度先验恢复 & 中、重两档各 3 次 & 满足磁观测前提时可复现闭环恢复 & 受确定性先验与窗口判定限制 \\
专用仓库算法级电缆边界 & 多场景确定性单次扫描（\(n=1\)） & 在线先验修正与自适应之字形为关闭即失败机制；投影连续与物理配准解耦；调优后之字形埋深进入 0.15 m 参考线 & 纯仿真单次证据，不等同主仓端到端、多种子或实物；30 m 为扫描下限非失效阈值（详见 \S\ref{sec:app-a-9}） \\
\end{longtable}
\endgroup

正文中对“旧口径被替代”的说明，应理解为分析定义或实验公平性得到修正，而不是删除不利结果。任何单次结果仍须标注 \(n=1\)，任何代理场景仍须与原生声磁耦合场景区分。

在上述压缩证据矩阵之外，本文另保留一份逐类实验的完整证据清单（表~\ref{tab:ch05-evidence-inventory}），供写作时逐条核查引用边界之用。其作用不是实验流水账，而是要求每条写入正文的结论都能在此找到对应的样本量与测量边界，从而避免在引用少数种子的实验数据时无意越界；各行末尾的交叉引用指向其正文主结果所在小节。

\begin{longtable}[]{@{}
  >{\raggedright\arraybackslash}p{(\linewidth - 6\tabcolsep) * \real{0.2308}}
  >{\raggedleft\arraybackslash}p{(\linewidth - 6\tabcolsep) * \real{0.3077}}
  >{\raggedright\arraybackslash}p{(\linewidth - 6\tabcolsep) * \real{0.2308}}
  >{\raggedright\arraybackslash}p{(\linewidth - 6\tabcolsep) * \real{0.2308}}@{}}
\caption{已完成实验及证据边界汇总}\label{tab:ch05-evidence-inventory}\tabularnewline
\toprule\noalign{}
\begin{minipage}[b]{\linewidth}\raggedright
实验类别
\end{minipage} & \begin{minipage}[b]{\linewidth}\raggedleft
样本量
\end{minipage} & \begin{minipage}[b]{\linewidth}\raggedright
可写结论
\end{minipage} & \begin{minipage}[b]{\linewidth}\raggedright
边界
\end{minipage} \\
\midrule\noalign{}
\endfirsthead
\toprule\noalign{}
\begin{minipage}[b]{\linewidth}\raggedright
实验类别
\end{minipage} & \begin{minipage}[b]{\linewidth}\raggedleft
样本量
\end{minipage} & \begin{minipage}[b]{\linewidth}\raggedright
可写结论
\end{minipage} & \begin{minipage}[b]{\linewidth}\raggedright
边界
\end{minipage} \\
\midrule\noalign{}
\endhead
\bottomrule\noalign{}
\endlastfoot
基线定位（多种子） & \(n=3\) & 基线定位指标均值\allowbreak{}±标准差可用 & 仅限 30 s 片段 \\
DVL 丢包四档（多种子） & \(n=3\)/档 & 四档均三次全部成功，支撑第 3 章鲁棒性 & 仅限定位指标 \\
磁畸变/声呐/综合四场景（多种子） & \(n=3\)/场景 & 四场景均三次全部成功，支撑第 3 章鲁棒性 & 水平面 RMSE 非单调 \\
不确定性感知控制主消融（定位） & 3 场景×2 模式×3 种子 & 基线与综合压力下定位改善，重度丢包无优势 & 离线滤波指标 \\
主消融控制侧指标 & 3 场景×2 模式×3 种子 & 横向 RMSE、控制变化率、回退率与非零求解时间均可观测 & 历史失败帧耗时字段不可与当前口径混用 \\
UA-MPC 机制消融与迭代上限对照 & P-1/P-2/H4 多档对照 & sigmoid 控制代价降权是可解性与横向精度主贡献 & 单纯提高迭代上限会阻塞控制周期 \\
分源 NIS 审计 & 8 场景×3 种子，35243 次量测更新 & 深度与 DVL 一致性偏差方向相反（见 §\ref{sec:ch05-5-5-5}） & 历史混合阈值只作操作性触发 \\
ES-EKF 五分钟时序敏感性 & 3 场景×5 种子×2 相位，30 次运行 & 成对噪声下隔离 0/100 ms DVL 相位效应（见 §\ref{sec:ch05-5-5-5}） & 离线动力学；无持续绝对水平位置观测 \\
实测背景噪声回放 & 同协方差高斯/实测回放×5 种子 & 均值接近，实测时序显著增大跨种子离散（见 §\ref{sec:ch05-5-5-5}） & 3.92 s 记录循环，不支持长时平稳性 \\
滤波策略与时间戳基线 & 3 场景×6 策略×5 种子，90 次运行 & Huber 更新显著抑制脉冲异常；分源 R 避免跨源混合（见 §\ref{sec:ch05-5-5-5}） & 离线动力学；时间戳膨胀不等同延迟状态更新 \\
基线定位/控制/决策 & \(n=1\) & 工具链与单次基准可用 & 不能替代多种子 \\
地形跟随 PID/MPC & \(n=1\)/组 & PID 地形跟随是当前近底主结果 & 需重复统计 \\
PID 地形跟随三档（多种子复验） & \(n=3\)/档 & 三档均三次全部成功、零安全违规、离底 RMSE 0.59--0.71 m（见 §\ref{sec:ch05-5-5-3}） & 部分档合并了重跑种子 \\
深度 MPC 调参 & 多轮调参 & 深度 MPC 不作为主线 & 只能写回退理由 \\
MPC 极端平面路径 & \(n=1\)+\(n=5\) 种子/场景 & 公平口径下 MPC 在多数极端路径优于或持平基线，受控扰动五种子统计复核（见 §\ref{sec:ch05-5-5-4}） & 仅直角路径 LOS 更优（诚实边界） \\
TMR8637 随附测试报告 & 1 份报告/编号 6\# 模组三轴 & 1 Hz 为 200 pT/\(\sqrt{\text{Hz}}\)，50 Hz 约 20--30 pT/\(\sqrt{\text{Hz}}\)；模组本体具备窄带噪声裕度（见 §\ref{sec:ch05-5-3-2}） & 供应方随附单件报告，非本文独立复测，非完整采集链 \\
SK2301 短接 ENOB 与 45 Hz 锁相 & 4 档采样率×10 s，最高 160000 样本 & ENOB 16.877--18.342 bit；45 Hz 矢量 RMS 0.0230 nT（见 §\ref{sec:ch05-5-3-2}） & ADC 子链灵敏度，不等同整机绝对精度 \\
ADC--TMR 全链路宿舍背景回放 & 2 段默认背景记录，共 13.104 s；17 段 raw data 全量审计 & 45 Hz 向量锁相中位 2.584 nT、50 Hz 工频中位 169.284 nT；45 Hz Biot--Savart 注入回放幅值 RMSE 0.792/0.271 nT（见 §\ref{sec:ch05-5-3-2}） & 非屏蔽宿舍环境，不能替代实验室本底计量；用于半实物噪声回放与压力测试 \\
ADC--TMR 实测背景 ROS2 闭环回放 & 清洁/同协方差高斯/实测回放三臂各 20 s & 噪声只进入磁观测量测副本；实测回放横偏 p95 为 7.767 m、平均置信度 0.8477、控制变化率 RMS 0.06346，未改变同配置闭环主导状态（见 §\ref{sec:ch05-5-3-2}） & Direction A 解耦轻量闭环门禁，非完整 AUV 水下实测；由于大横偏和埋深方差缺失，不作为行业验收通过率实验 \\
45 Hz 电缆--TMR--ArUco 联合反演 & 106821 ADC 样本，777 对齐点 & 复 I/Q 自由/固定电流拟合 \(R^2=0.8792/0.8711\)（见 §\ref{sec:ch05-5-4}） & 直导线与公共复尺度模型 \\
磁传感器杆臂/安装角标定（仿真） & \(n=1\) 全流程 & 杆臂改正链路与文件契约成立，平移降 96.27\%、旋转降 95.31\%、通过（见 §\ref{sec:ch05-5-3-3}） & 仿真标定框架、非真机转台、未估九参数 \\
代理电缆六场景正式矩阵 & 6 场景×2 模式×3 种子 & 36/36 完成、回退为 0；保守 UA 将全局控制变化率 RMS 由 0.247 降至 0.085（见 \S\ref{sec:ch05-5-7-7}） & 横向 RMSE 未系统改善；代理配置不代表原生声磁闭环 \\
CBF 安全过滤低/中/高三种子矩阵 & 低/中/高 \(\times\) 3 种子 \(\times\) PID/MPC，各约 60 s & 18/18 运行有效；最小净空 2.70--2.90 m，1.5 m 低净空违规率和穿底率均为 0；MPC 回退为 0，求解耗时 p95 最大 23.84 ms（见 \S\ref{sec:ch05-5-7-7}） & 手动设定点加运动学代理口径；不代表 PVS 原生执行链或所有极端场景 \\
原生电缆因子生成器 & 16 组合×2 种子，32 次契约验证 & 几何、地形、实测噪声与横流可独立开关，真值与量测隔离 & 生成器验证，不替代控制闭环因子交互实验 \\
R22 原生因子闭环全矩阵 & \(2^4\) 组合×3 种子×2 模式，96 次 20 s & 96/96 有效，回退与安全违规均为 0；最坏横向 p95 为 1.716 m，最小净空 2.536 m；求解耗时 p95 最大 24.02 ms（见 \S\ref{sec:ch05-5-7-7}） & 离线仿真代理口径；非 PVS 原生执行链或真机物理时延 \\
R22 60 s 全因子长时复核 & \(2^4\) 组合×3 种子×2 模式，96 次 60 s & 96/96 有效，回退与安全违规均为 0；最坏横向 p95 为 2.176 m，最小净空 2.384 m；求解耗时 p95 最大 14.93 ms（见 \S\ref{sec:ch05-5-7-7}） & 未触发路径终点；仍非 PVS 原生执行链/真机 \\
R22 地形加横流 120 s 终端边界测试 & 4 组合×3 种子×2 模式，24 次 120 s & 24/24 有效，安全违规为 0；全窗口最坏横向 p95 为 6.839 m；终端审计显示活动巡缆段 p95 最大 2.152 m、终端尾段 p95 最大 7.703 m；UA 回退最大 0.113（见 \S\ref{sec:ch05-5-7-7}） & 终端任务边界；不能写成活动段跟踪崩溃或长航时定点保持 \\
R22 120 s 任务层终端退出验证 & 同上，到达路径末端即退出 & 24/24 有效且均到达路径末端；退出后最坏横向 p95 为 2.152 m，UA 回退率最大 0.018；solver p95/max 最大 22.99/53.27 ms（见 \S\ref{sec:ch05-5-7-7}） & 消除终端尾段指标污染；不代表定点保持或 PVS 原生执行链 \\
PVS 原生执行链冒烟 & 约 50 s 定深定向仿真 & 审计通过；PVS 原生定深定向自驾配置，AUTO 日志 149 次，运动学日志为 0；设定点/控制/状态帧分别 245/720/998，solver p95/max 为 19.32/20.40 ms（见 \S\ref{sec:ch05-5-7-7}） & 执行链证据；非 R22 性能矩阵、真机或海试 \\
PC104 UDP 时序探针 & 30 s 主机中继安全零推力实机记录 & 下行 300 帧进入中继，解析 450 帧 \texttt{\$AUV} 上行；到达频率约 15.0 Hz，p50/p95/p99 间隔约为 58.0/85.7/85.9 ms，解析错误为 0 & 只报告到达间隔/抖动；无共享时钟或固件回显时不作单向物理延迟结论 \\
行为树 vs 状态机 & 1000 次延迟、100 条噪声序列、500 次组合故障 & 等价规则下响应和生存率相同；迟滞去抖降低颤振 & 结论限于决策核心仿真 \\
海缆行业标准数字孪生验收 & 3 次全新运行 & 全链路闭环、三次全部就绪且通过（见 §\ref{sec:ch05-5-5-10}） & 数字孪生确定性先验，非真实检测噪声 \\
解耦轻量闭环 & 单次冒烟验证（\(n=1\)） & 满足磁观测前提时在线先验修正被真实观测接受，算法实机部署接口成立（见 §\ref{sec:ch05-5-5-11}） & 轻量运动学闭环，无地磁背景/检测噪声/硬件时延/六自由度水动力 \\
六自由度闭环偏差先验恢复 & 中/重档各 3 次全新运行 & 在线修正被真实接受、六自由度闭环恢复得到复现，中/重档各三次全部就绪通过（见 §\ref{sec:ch05-5-5-11}） & 数字孪生确定性先验、静态位姿扭曲、缆在车下满足前提、窗口内判定、每档三种子 \\
\end{longtable}

\section{海缆巡检验收判据}\label{sec:app-a-8}

数字孪生验收采用起点健康、有效窗口、逐帧质量和多次运行聚合四级判据。令 \(s_k\) 为路由进度，\(d_{\perp,k}\) 为横向偏移，\(q_k\in\{0,1\}\) 表示埋深估计是否就绪，\(I_k\in\{0,1\}\) 表示第 \(k\) 帧是否进入验收窗口。有效窗口定义为
\[
  I_k=
  \mathbb{I}
  \left(
    q_k=1,\;
    s_k\le50\,\mathrm{m},\;
    |d_{\perp,k}|\le2.0\,\mathrm{m}
  \right).
\]
起点健康检查使用前 \(30\) 个样本，并要求起始路由进度不超过 \(20\,\mathrm{m}\)、起始横向偏移不超过 \(5\,\mathrm{m}\)。该宽松起点阈值用于排除明显错误初始化，不代替正式验收走廊。

令 \(\sigma_{d,k}\) 为埋深估计标准差，\(\sigma_{\max}\) 为规定上限，则不确定度超限比例为
\[
  r_\sigma=
  \frac{
    \sum_k I_k\,
    \mathbb{I}(\sigma_{d,k}>\sigma_{\max})
  }{
    \sum_k I_k
  }.
\]
本文要求 \(r_\sigma\le0.05\)。多次运行聚合至少包含 \(N_{\min}=3\) 次独立运行，通过率定义为
\[
  r_{\mathrm{pass}}=
  \frac{N_{\mathrm{pass}}}{N_{\mathrm{run}}},
  \qquad
  r_{\mathrm{pass}}\ge0.67.
\]

三次数字孪生运行均满足就绪与通过判据，但“初步验收就绪”只表示数据完整性、窗口有效性和工程阈值在数字孪生中一致，不表示已通过真实海缆检测精度验收。

\section{端到端证据层级与恢复条件}\label{sec:app-a-9}

端到端电缆探测由磁场、导航和声学观测共同驱动，每个时刻输出至少包括横向偏移 \(d_\perp\)、路由进度 \(s\)、埋深估计 \(\hat d\)、埋深标准差 \(\sigma_d\)、置信度 \(c\)、磁信噪比和质量状态。具体程序接口、字段名称和消息通道属于软件文档，不影响本文算法定义，故不在此列出。

设先验电缆点为 \(\boldsymbol p_i\)，受控先验畸变写为
\[
  \tilde{\boldsymbol p}_i
  =
  \boldsymbol R(\delta\psi)\boldsymbol p_i
  +\boldsymbol t,
\]
其中 \(\delta\psi\) 为平面旋转偏差，\(\boldsymbol t\) 为平移偏差。在线修正估计 \(\hat{\delta\psi}\) 与 \(\hat{\boldsymbol t}\)，并仅在磁观测满足质量门限时接受更新。该设计保证“恢复”来自新观测对先验的校正，而不是直接放宽验收阈值。

端到端证据按以下层级解释：

\begin{enumerate}
  \item \textbf{算法级扫描：}用于识别同源算法对先验偏差、曲率和观测失效的敏感性，属于纯仿真单次证据。
  \item \textbf{开环回放：}证明先验畸变会沿数据链确定性传导至横向偏移和验收失败，但不包含车辆运动反馈。
  \item \textbf{解耦轻量闭环：}证明观测接受、在线修正和指导量更新可以形成闭环，但运动模型经过简化。
  \item \textbf{六自由度闭环：}在满足磁观测物理前提时，中、重两档各三次复现窗口内恢复。
  \item \textbf{实物验收：}TMR8637 已有单件模组三轴随附测试报告，完整声磁噪声、采集链、硬件标定、网络时延和水域试验仍未完成。
\end{enumerate}

六自由度恢复实验中，电缆相对车体的名义垂直间距约为 \(7.5\,\mathrm{m}\)，之字形横切幅度为 \(0.6\)，验收走廊半宽为 \(3.4\,\mathrm{m}\)。这些量是恢复实验的物理前提与判定尺度，不是通用海缆参数。重度先验畸变由约 \(10.25\,\mathrm{m}\) 的起始横偏在约 \(12\,\mathrm{s}\) 内收敛进入走廊；关闭在线修正的对照不能形成同等恢复。

\S\ref{sec:ch05-5-5-11} 引用的算法级结论均来自专用磁探测仓库的确定性离线扫描（审计基线提交 \texttt{12e1884b}），其证据契约在此补足，以便正文数字可逐项回溯。场景族的关键参数包括：巡检车道长度与 \(50/70\,\mathrm{m}\) 车道间距、U 形掉头半径、以及先验畸变的轻/中/重三档（对应上式中平移 \(\boldsymbol t\) 与旋转 \(\delta\psi\) 的取值梯度）。判据方面，单次运行的 \texttt{passed} 由该仓库的 DL/T 1278 风格复合门限给出，其中 \texttt{route\_progress\_large\_jump\_count} 定义为沿先验路由的单步进度增量超过跨车道判定阈值的次数，用于捕获错位先验在 U 弯处诱发的跨车道投影跳变。部署接口的两组量语义不同且不可混用：\texttt{map\_frame\_*} 描述投影落在先验哪一段的地图系连续性，即使任务已失败其单步跳变仍可保持在约 \(0.2\,\mathrm{m}\)；\texttt{prior\_alignment\_*} 描述把错位先验拉向真实电缆的物理平移/旋转配准量，是任务能否恢复的真正状态。故“投影连续”“物理配准”“任务完成”是三个不同概念，不能相互替代。曲率一项，\(30\,\mathrm{m}\) 是扫描环境的硬下限而非已击穿的失效阈值：在 \(30\text{--}120\,\mathrm{m}\) 范围内未观测到弯道触发的失效，瓶颈接近电缆几何物理下限。真值字段（如真实电缆最近点坐标）仅进入离线评价，不进入在线感知管线。上述扫描随机种子未透传、当前为 \(n=1\)，因此只能作为同源算法的鲁棒性边界，不能改写为主仓端到端、多种子或实物结论。

该算法级证据的最小文件集、结果 CSV、重绘脚本与迁移到本文的方法图/因果图，已固化为证据清单条目 \texttt{MAG-CABLE-ALGORITHM-BOUNDARY}（数据层 \texttt{algorithm\_simulation}、等级 C、状态 \texttt{complete\_with\_boundaries}），其中每个文件均带 SHA256 校验值，可据以逐项回溯正文引用的算法级数字。

\section{代理极端场景的因子化定义}\label{sec:app-a-10}

原生仿真后端尚不能完整耦合电缆几何、近底地形、声磁观测和横流。为验证控制闭环在组合压力下能否运行，本文将代理场景表示为四个因子的组合
\[
  \mathcal{S}
  =
  \mathcal{G}
  \times
  \mathcal{T}
  \times
  \mathcal{O}
  \times
  \mathcal{C},
\]
其中 \(\mathcal{G}\) 表示电缆中心线几何，\(\mathcal{T}\) 表示地形坡度与起伏，\(\mathcal{O}\) 表示声磁观测质量，\(\mathcal{C}\) 表示海流速度和方向。

\begingroup
\small
\setlength{\tabcolsep}{3pt}
\renewcommand{\arraystretch}{1.04}
\begin{longtable}{@{}p{0.23\linewidth}p{0.27\linewidth}p{0.40\linewidth}@{}}
\caption{代理极端场景的主导因子}\label{tab:app-a-proxy-scenarios}\\
\toprule
场景 & 主导因子 & 主要考察问题 \\
\midrule
\endfirsthead
\toprule
场景 & 主导因子 & 主要考察问题 \\
\midrule
\endhead
\bottomrule
\endfoot
连续 S 弯 & \(\mathcal{G}\) & 连续曲率变化下的预瞄与横向跟踪 \\
发卡掉头 & \(\mathcal{G}\) & 大航向变化和终端任务定义 \\
坡面穿越 & \(\mathcal{T}\) & 平面制导与近底高度控制耦合 \\
埋设间断 & \(\mathcal{O}\) & 观测中断下的置信度与回退 \\
横流冲击 & \(\mathcal{C}\) & 外部扰动下的横向保持能力 \\
综合极端 & \(\mathcal{G},\mathcal{T},\mathcal{O},\mathcal{C}\) & 多因素同时作用时的闭环可运行性 \\
\end{longtable}
\endgroup

六类场景分别在固定权重和不确定性感知两种模式下完成单种子冒烟测试，共 \(6\times2=12\) 次运行。该结果只能说明代理压力测试链路能够生成有效控制指标；由于每个组合仅 \(n=1\)，不能据此宣称两种控制模式存在统计显著差异，也不能替代原生声磁耦合数字孪生或海试。

\section{近底安全与原生因子闭环的补充矩阵}\label{sec:app-a-11}

本节汇总正文 §\ref{sec:ch05-5-7-7} 只给出结论的两张多种子矩阵，供逐条复核。表~\ref{tab:ch05-cbf-terrain-matrix}给出控制屏障函数（CBF）过滤器在低、中、高三档地形与 PID/MPC 两类控制器下的 18 个约 60 s 运行结果；正文已据此写明 1.5 m 安全违规率与海床穿透率均为 0、最小净空下界 2.70--2.90 m、MPC 求解耗时 p95 上界 23.84 ms 等边界结论。

\begin{longtable}[]{@{}
  >{\raggedright\arraybackslash}p{0.13\linewidth}
  >{\centering\arraybackslash}p{0.08\linewidth}
  >{\centering\arraybackslash}p{0.11\linewidth}
  >{\centering\arraybackslash}p{0.11\linewidth}
  >{\centering\arraybackslash}p{0.12\linewidth}
  >{\centering\arraybackslash}p{0.13\linewidth}
  >{\centering\arraybackslash}p{0.10\linewidth}@{}}
\caption{CBF 地形低/中/高 60 s 三种子复核矩阵结果}\label{tab:ch05-cbf-terrain-matrix}\tabularnewline
\toprule\noalign{}
运行 & 成功次数 & 最小净空下界（m） & 平均净空均值（m） & 违规/穿底最大值 & CBF 激活率/最低缩放 & 求解耗时 p95 最大（ms） \\
\midrule\noalign{}
\endfirsthead
\toprule\noalign{}
运行 & 成功次数 & 最小净空下界（m） & 平均净空均值（m） & 违规/穿底最大值 & CBF 激活率/最低缩放 & 求解耗时 p95 最大（ms） \\
\midrule\noalign{}
\endhead
\bottomrule\noalign{}
\endlastfoot
PID·低 & 3/3 & 2.90 & 3.19 & 0 / 0 & 0.669 / 0.678 & N/A \\
PID·中 & 3/3 & 2.90 & 3.14 & 0 / 0 & 0.473 / 0.563 & N/A \\
PID·高 & 3/3 & 2.80 & 3.16 & 0 / 0 & 0.330 / 0.413 & N/A \\
MPC·低 & 3/3 & 2.90 & 3.22 & 0 / 0 & 0.775 / 0.585 & 23.84 \\
MPC·中 & 3/3 & 2.80 & 3.15 & 0 / 0 & 0.544 / 0.500 & 19.62 \\
MPC·高 & 3/3 & 2.70 & 3.21 & 0 / 0 & 0.463 / 0.180 & 23.62 \\
\end{longtable}

表~\ref{tab:ch05-r22-native-factorial}给出 R22 原生因子闭环全矩阵与时序审计。该矩阵覆盖曲线几何、地形/埋设、实测磁噪声回放与横流四个因子的 \(2^4\) 个组合，每组 3 个种子，并分别运行固定权重基线与保守不确定性感知控制，共 96 次、每次 20 s；正文已据此写明 96/96 有效、回退与安全违规均为 0、最坏横向 p95 为 1.716 m、最小净空 2.536 m、求解耗时 p95/max 上界 24.02/38.83 ms 且相对三级周期预算未越界等结论。

\begingroup
\small
\setlength{\tabcolsep}{3pt}
\begin{longtable}[]{@{}
  >{\raggedright\arraybackslash}p{0.16\linewidth}
  >{\centering\arraybackslash}p{0.12\linewidth}
  >{\centering\arraybackslash}p{0.13\linewidth}
  >{\centering\arraybackslash}p{0.12\linewidth}
  >{\centering\arraybackslash}p{0.16\linewidth}
  >{\centering\arraybackslash}p{0.10\linewidth}@{}}
\caption{R22 原生因子闭环全矩阵与时序审计}\label{tab:ch05-r22-native-factorial}\tabularnewline
\toprule\noalign{}
运行口径 & 有效运行 & 最坏横向 p95 & 最小净空 & 求解耗时 p95 / max & 周期越界 \\
\midrule\noalign{}
\endfirsthead
\toprule\noalign{}
运行口径 & 有效运行 & 最坏横向 p95 & 最小净空 & 求解耗时 p95 / max & 周期越界 \\
\midrule\noalign{}
\endhead
\bottomrule\noalign{}
\endlastfoot
\(2^4\)×3 种子×2 模式，20 s & 96/96 & 1.716 m & 2.536 m & 24.02 / 38.83 ms & 0 \\
\end{longtable}
\endgroup

\section{证据边界、负结果归因与剩余验证}\label{sec:app-a-12}

本节把全章实验的负结果、口径限制和不可外推项集中收束为问题域，而不是逐条复述前文结果。表~\ref{tab:ch05-negative-results}保留各类边界的原因、处置和剩余验证指向；正文各结果小节只保留正向主结果与最小必要限定，相同局限统一指向本表。需要说明的是，本表``剩余验证''列并非未知缺口清单，而是已按可独立交付的实物证据链拆解的推进路径——每项均可在当前仿真结论之上增量验证，其向实物层推进的路线见 §\ref{sec:ch06-6-3}，不构成对已有结论的否定。

\begingroup
\small
\setlength{\tabcolsep}{3pt}
\renewcommand{\arraystretch}{1.08}
\begin{longtable}[]{@{}p{0.16\linewidth}p{0.28\linewidth}p{0.31\linewidth}p{0.20\linewidth}@{}}
\caption{证据边界与剩余验证汇总}\label{tab:ch05-negative-results}\tabularnewline
\toprule\noalign{}
问题域 & 边界归因 & 已采取处置与可写结论 & 剩余验证 \\
\midrule\noalign{}
\endfirsthead
\toprule\noalign{}
问题域 & 边界归因 & 已采取处置与可写结论 & 剩余验证 \\
\midrule\noalign{}
\endhead
\bottomrule\noalign{}
\endlastfoot
磁探测计量与验收口径 & 45 Hz 短背景高于 0.05 nT，受矩形窗泄漏、三轴相关与短时结构影响；复 I/Q 与埋深联合拟合仍基于直导线和公共复尺度；数字孪生验收限定于干净先验、静态位姿扭曲和有效巡检窗口 & 已记录窗函数、相干增益和 ENBW，完成高斯/实测回放五种子、复 I/Q 拟合、合成可辨识性和 ENOB 测算；可写为数字孪生窗口内三次就绪通过、中重档各三次闭环验收通过，不写成现场行业验收或系统级 0.05 nT 计量结论 & 长时 M0/M1/M4 重复采集、可溯源电流/距离/转台、九参数标定、完整不确定度预算和真实检测噪声下实物验收 \\
ES-EKF 统计一致性 & 历史实现把一维深度 NIS 与三维 DVL NIS 混入同一窗口；长航段缺持续绝对水平位置/航向约束，时间差膨胀 \(R\) 不能消除相位与积分漂移 & 固定/全局/分源/Huber 对照与分源一致性 A/B 已完成，深度回到卡方带内；两相位成对对照确认水平漂移不是单靠时间差膨胀可解 & 下调 DVL 名义协方差并连同主线默认重跑矩阵；加入延迟状态更新与声磁绝对水平观测闭环 \\
在线先验与几何前提 & 初次在线先验修正失败来自产磁电缆与车近共面，违反直线埋缆横偏反演前提；该失败是残差门正确拒绝，而非部署接口失效 & 迁回满足前提的几何后，约 98\% 观测被接受，中、重档各三次通过；可写为在线修正接口与闭环恢复机制成立，前提是满足磁观测几何 & 真实背景、动态先验误差、硬件时延与非理想几何下复验 \\
UA/R13-v2 控制证据 & 旧代理闭环高回退来自硬俯仰饱和与低频快照口径；保守分源 UA 主要改变控制风格，未系统改善横向 RMSE；综合极端因 \(p_{\mathrm{track}}\) 未越过进入门而不进入跟踪态，原 36 次未落盘速度/授权遥测 & 有界可导饱和后 36 快照全收敛、18 组关键闭环零回退；主结论收束为控制平滑与实时可解性（控制变化率 RMS 0.247→0.085），RMSE 边界由表~\ref{tab:ch05-r13-v2-main-results} 给出；补 12 次遥测证明速度缩放，但不回写正式矩阵结论 & 重构分源置信接口并重跑主矩阵；提升综合极端观测质量；正式矩阵内落盘速度门控与授权遥测，并在真机口径下复验精度效应 \\
CBF/R22 与 PVS 原生闭环 & CBF 矩阵仍是手动设定点加运动学代理；R22 全矩阵为离线代理，PVS 原生冒烟仅覆盖定深定向；120 s 固定时长会把参考路径完成后的尾段污染进巡缆指标；CBF/R22 p95 上界约 24 ms，不等于 10 ms 硬实时 & CBF 18/18 有效，最小净空 2.70--2.90 m，违规和穿底为 0；R22 96/96 有效、回退和安全违规为 0，活动段 p95 最大 2.152 m，任务层到达即退出后 24/24 退出；相对 20 Hz/10 Hz/0.3 s 三级周期预算未越界 & PVS 原生全因子安全/性能矩阵、速度/横流/初始偏差边界搜索、终端保持/横流悬停控制、长航时稳定闭环与真机近底复验 \\
实物通信与端到端时延 & PC104 探针仅能报告主机中继 30 s 内的到达间隔和抖动，缺共享时钟或固件回显，不能反推出物理单向时延 & 已记录 450 帧上行、约 15 Hz 到达、解析错误为 0；该结果只作为序列完整性与软件链路冒烟审计 & 时钟同步或固件回显下测量单向/往返物理时延，并完成 Jetson--PC104 端到端闭环 \\
\end{longtable}
\endgroup

这些结果共同表明，当前方法的主要算法链已经闭合，剩余工作集中在统计标定、分源置信映射、原生耦合闭环、目标平台实时复验和系统级计量验证五个层面。它们限定结论的适用范围，但不改变本文关于低成本窄带采集、声磁协同估计、不确定性感知控制和闭环恢复机制的核心论证。
